\chapter{Lesson 3: Greetings and Names}

\section{Vocabulary}

\begin{center}
\begin{tabular}{llll}
\toprule
Characters & Pinyin & Part of speech & Meaning \\
\midrule
\zh{你} & nǐ & pronoun & you \\
\zh{您} & nín & pronoun & you (polite) \\
\zh{我} & wǒ & pronoun & I, me \\
\zh{他} / \zh{她} & tā & pronoun & he / she \\
\zh{们} & men & suffix & plural for people (\zh{我们} we, \zh{你们} you all) \\
\zh{好} & hǎo & adjective & good, well \\
\zh{请问} & qǐngwèn & verb & may I ask \\
\zh{贵姓} & guìxìng & noun & (your) honourable surname \\
\zh{姓} & xìng & verb / noun & to be surnamed; surname \\
\zh{叫} & jiào & verb & to be called \\
\zh{什么} & shénme & question pronoun & what \\
\zh{名字} & míngzi & noun & name \\
\zh{是} & shì & verb & to be \\
\zh{不} & bù & adverb & not, no \\
\zh{吗} & ma & particle & yes/no question particle \\
\zh{呢} & ne & particle & ``and \ldots?'' (follow-up question) \\
\zh{也} & yě & adverb & also, too \\
\zh{老师} & lǎoshī & noun & teacher \\
\zh{学生} & xuésheng & noun & student \\
\zh{先生} & xiānsheng & noun & Mr., husband, sir \\
\zh{小姐} & xiǎojiě & noun & Miss, young lady \\
\zh{中国人} & Zhōngguórén & noun & Chinese person \\
\zh{加拿大人} & Jiānádàrén & noun & Canadian \\
\zh{美国人} & Měiguórén & noun & American \\
\zh{哪国人} & nǎ guó rén & phrase & which nationality \\
\bottomrule
\end{tabular}
\end{center}

\section{Greetings}

\begin{itemize}
  \item \zh{你好！} \emph{Nǐ hǎo!} is the all-purpose ``hello''. To an older person or a
        teacher, \zh{您好} \emph{nín hǎo} is more respectful.
  \item \zh{你好吗？} \emph{Nǐ hǎo ma?} ``How are you?'' is used less than English speakers
        expect; it implies you haven't seen the person for a while. Answer with
        \zh{我很好，谢谢。} \emph{Wǒ hěn hǎo, xièxie.}
  \item Among friends, Chinese speakers often greet by asking what the other person is
        doing: \zh{吃了吗？} \emph{Chī le ma?} ``Have you eaten?'', or \zh{去上课啊？}
        ``Off to class?''. These are greetings, not real questions.
\end{itemize}

\begin{aside}
\zh{您} \emph{nín} contains \zh{心} ``heart'' under \zh{你}, a nice reminder that it is
the polite form. It has no plural: ``\zh{您们}'' is avoided in speech; say
\zh{你们} or \zh{各位} \emph{gèwèi} ``everyone''.
\end{aside}

\section{Asking and giving names}

\subsection{Surname: \zh{姓}}

\begin{example}[Polite exchange]\leavevmode
\begin{itemize}
  \item \textbf{A:} \zh{请问，您贵姓？} \emph{Qǐngwèn, nín guìxìng?} May I ask your surname?
  \item \textbf{B:} \zh{我姓李。你呢？} \emph{Wǒ xìng Lǐ. Nǐ ne?} My surname is Li. And you?
  \item \textbf{A:} \zh{我姓陈，叫陈明。} \emph{Wǒ xìng Chén, jiào Chén Míng.} My surname is Chen;
        my name is Chen Ming.
\end{itemize}
\end{example}

\begin{important}[\zh{姓} is a verb]
\zh{我姓王。} ``I am surnamed Wang.'' Do \emph{not} say *\zh{我是姓王} or answer
\zh{贵姓} with *\zh{我贵姓王}: \zh{贵} is an honorific for the other person and is never
used about yourself. Negation: \zh{我不姓王。}
\end{important}

\subsection{Full name: \zh{叫}}

\zh{叫} is followed by a full name or a given name, never by a surname alone.
\begin{itemize}
  \item \zh{你叫什么名字？} \emph{Nǐ jiào shénme míngzi?} What's your name?
  \item \zh{我叫王小文。} \emph{Wǒ jiào Wáng Xiǎowén.} I'm called Wang Xiaowen.
  \item \zh{他叫什么？} \emph{Tā jiào shénme?} What is he called?
\end{itemize}

\begin{aside}
Chinese names put the surname first. Most surnames are one syllable (\zh{王 李 张 刘 陈});
a few are two (\zh{欧阳} Ōuyáng, \zh{司马} Sīmǎ). Given names are one or two syllables
and are chosen for their meaning. A Chinese person with an English name often writes
it as ``Xiaowen (Wendy) Wang''. Among friends, \zh{小} + surname (\zh{小王}) is used for
someone young and \zh{老} + surname (\zh{老王}) for someone older; titles follow the
name: \zh{王老师}, \zh{李先生}, \zh{陈小姐}.
\end{aside}

\section{Grammar}

\subsection{\zh{是} sentences}

\zh{是} links two nouns: A \zh{是} B ``A is B''. It is negated with \zh{不}
(pronounced \emph{bú} before \emph{shì}).
\begin{itemize}
  \item \zh{我是学生。} \emph{Wǒ shì xuésheng.} I am a student.
  \item \zh{她不是老师。} \emph{Tā bú shì lǎoshī.} She is not a teacher.
  \item \zh{李先生是加拿大人。} Mr.\ Li is Canadian.
\end{itemize}

\begin{important}[No \zh{是} before adjectives]
Adjectives act as predicates directly, usually with \zh{很} in front:
\zh{我很好}, not *\zh{我是好}. The \zh{很} is often unstressed and does not really mean
``very''; without it, the sentence sounds contrastive (\zh{我好，他不好}).
\end{important}

\subsection{Yes/no questions with \zh{吗}}

Add \zh{吗} to the end of a statement. The word order does not change.

\begin{center}
\begin{tikzpicture}[w/.style={draw, rounded corners, fill=blue!8, minimum height=0.8cm, minimum width=1.3cm}]
  \node[w] (s) at (0,0) {\zh{你}};
  \node[w] (v) at (1.6,0) {\zh{是}};
  \node[w] (o) at (3.4,0) {\zh{中国人}};
  \node[w, fill=red!12] (q) at (5.3,0) {\zh{吗}};
  \node[font=\small] at (6.2,0) {?};
  \node[font=\scriptsize, below=0.1cm of s] {subject};
  \node[font=\scriptsize, below=0.1cm of v] {verb};
  \node[font=\scriptsize, below=0.1cm of o] {object};
  \node[font=\scriptsize, below=0.1cm of q] {particle};
  \node[w, fill=gray!10] (s2) at (0,1.5) {\zh{你}};
  \node[w, fill=gray!10] (v2) at (1.6,1.5) {\zh{是}};
  \node[w, fill=gray!10] (o2) at (3.4,1.5) {\zh{中国人}};
  \node[font=\small, anchor=east] at (-0.8,1.5) {statement:};
  \node[font=\small, anchor=east] at (-0.8,0) {question:};
  \draw[->, gray] (o2) -- (o);
  \node[font=\scriptsize, red!70!black] (lab) at (5.3,1.5) {add \zh{吗}};
  \draw[->, thick, red!70!black] (lab) -- (q.north);
\end{tikzpicture}
\end{center}

Answers repeat the verb rather than saying ``yes'' or ``no'':
\begin{itemize}
  \item \zh{你是中国人吗？}---\zh{是，我是中国人。} / \zh{不是，我是加拿大人。}
\end{itemize}
\zh{对} \emph{duì} ``correct'' can also be used to confirm.

\subsection{Question-word questions}

Put the question word in the position of the answer; do not move it to the front and do
not add \zh{吗}.
\begin{itemize}
  \item \zh{你叫什么名字？}---\zh{我叫\textbf{王朋}。}
  \item \zh{他是哪国人？}---\zh{他是\textbf{法国}人。} What nationality is he?
  \item \zh{她是谁？}---\zh{她是\textbf{我的老师}。} Who is she?
\end{itemize}

\begin{important}[No \zh{吗} with question words]
*\zh{你叫什么名字吗？} is wrong: the question word already makes it a question.
\end{important}

\subsection{Follow-up questions with \zh{呢}}

After a statement, ``X \zh{呢}？'' asks the same thing about X.
\begin{itemize}
  \item \zh{我是学生，你呢？} I'm a student; and you?
  \item \zh{我姓张。他呢？}---\zh{他也姓张。}
\end{itemize}

\subsection{The adverb \zh{也}}

\zh{也} ``also'' must come \emph{after} the subject and \emph{before} the verb. It
never begins a sentence.
\begin{itemize}
  \item \zh{我也是学生。} I am also a student. (not *\zh{也我是学生})
  \item \zh{他也不是老师。} He isn't a teacher either.
\end{itemize}

\section{Model conversation}

At a Renison welcome event, two first-years meet.
\begin{itemize}
  \item \textbf{A:} \zh{你好！我叫安娜。你叫什么名字？}
  \item \textbf{B:} \zh{你好，安娜。我叫高明。}
  \item \textbf{A:} \zh{你是中国人吗？}
  \item \textbf{B:} \zh{不是，我是加拿大人。我爸爸妈妈是中国人。你呢？}
  \item \textbf{A:} \zh{我也是加拿大人。你是学生吗？}
  \item \textbf{B:} \zh{是，我是一年级的学生。}
\end{itemize}
(\zh{一年级} \emph{yī niánjí} = first year.)

\section*{Review}
\begin{itemize}
  \item \zh{您贵姓？}/\zh{我姓\ldots}: \zh{姓} is a verb; never use \zh{贵} about yourself.
  \item \zh{你叫什么名字？}/\zh{我叫\ldots}: \zh{叫} takes a full or given name.
  \item A \zh{是} B; negate with \zh{不是}; no \zh{是} before adjectives (\zh{我很好}).
  \item \zh{吗} turns a statement into a yes/no question; question words stay in place,
        with no \zh{吗}.
  \item X \zh{呢}？ asks ``and X?''; \zh{也} goes right before the verb.
  \item Chinese names: surname first; titles after the name.
\end{itemize}
